\documentclass[11pt, oneside]{article}   	% use "amsart" instead of "article" for AMSLaTeX format
\usepackage[margin=2cm]{geometry}
\geometry{letterpaper}
\usepackage[parfill]{parskip}    		% Activate to begin paragraphs with an empty line rather than an indent
\usepackage{graphicx}
\usepackage{amssymb}
\usepackage{amsmath}
\usepackage{multicol}
\usepackage{booktabs}
\usepackage{array}

%SetFonts


\title{uOttawa - Telfer \\ Workload Allocation Problem}
\author{JP, AS, TA, KV}
%\date{}							% Activate to display a given date or no date

\begin{document}
\maketitle


%%%%%%%%%%%%%%%%
%%% DESCRIPTION
%%%%%%%%%%%%%%%%
\section{Description}
\subsection{Intro}
The goal of this project is to create a course schedule. The schedule must
assign every course section to a professor, a semester, a time, and a
classroom. The implementation identifies a declared section by its course,
section number, semester, and time-of-day category. Its assignment also records
the professor, timeslot, and classroom. The notation $X_{pcistk}$ below suppresses
the declared time-of-day index for readability; all section-indexed constraints
are applied separately to each declared $(c,i,s,td)$ key.

The model has two kinds of requirements. Hard requirements describe what a
valid schedule must satisfy. Soft requirements describe desirable qualities,
such as respecting professor preferences or avoiding room changes. Because a
real instance may not satisfy every hard requirement, the solver handles the
requirements hierarchically: it first minimizes selected violation measures, then
optimizes the softer qualities within the resulting limits. A violation is not
proved unavoidable merely because a time-limited solve retains it.

Professor preferences are represented on a simple three-point numerical scale.
This keeps the input understandable and limits the amount of subjective
weighting introduced by the model. The scale can be refined later using
feedback from the people who use the resulting schedules.
%Faculty requirements naturally lend themselves to be constraints in the mathematical formulation. In addition to faculty requirements some implicit constraints need to be formulated. Faculty requirement constraints include section requirements for each course, professor qualifications, attendance and ca- pacity matching, etc. Implicit constraints include things like each professor can only teach once per time block, classroom can only be occupied once per time block, etc.


%%%%%%%%%%%%%%%%
%%% OPTIMIZATION MODEL
%%%%%%%%%%%%%%%%
\newpage
\section{Optimization Model}





%%%%%%%%%%%%%%%%
%%% SETS
%%%%%%%%%%%%%%%%
\subsection{Sets}
\begin{multicols}{2}
	$P$ - professors - "part-timer" is added to set\\
	$C$ - courses\\
	$I$ - unique section ID \\ 
	$K$ - classrooms - "other" classroom is added to set\\
	$S$ - semesters \\
	$T$ - timeslots  - "placeholder" is added for each $TD$\\
	$TD$ - times of day\\
	$DW$ - days of week \\
	$BL$ - blocks
\end{multicols}





%%%%%%%%%%%%%%%%
%%% SUBSETS
%%%%%%%%%%%%%%%%
\subsection{Subsets}
\begin{multicols}{2}
	$CS_{cis,td}$ - possible course $C$ section $I$ per $S$, $TD$ \\
	$TTD_{t,td}$ - timeslots to times of day \\
	$TDW_{t,dw}$ - timeslots to days of week \\
	$TBL_{t,bl}$ - timeslots to blocks
\end{multicols}

For readability the equations below use the derived index sets these subsets
induce, rather than restating the membership test inside every summation:

\begin{multicols}{2}
	$T_{td} = \{ t : [t,td] \in TTD \}$ \\
	$T_{dw} = \{ t : [t,dw] \in TDW \}$ \\
	$T_{bl} = \{ t : [t,bl] \in TBL \}$ \\
	$P^{*} = P \setminus \{ \text{part-timer} \}$
\end{multicols}





%%%%%%%%%%%%%%%%
%%% DATA PARAMETERS
%%%%%%%%%%%%%%%%
\subsection{Data Requirements}


\subsubsection{Time Data}
$tbw_{bl}$ - timeslots weights for block $BL$ (for day of week equations)  \\
$tc_{t,t}$ - timeslots concurrencies \\
$tb_{t,t}$ - timeslots that are back to back


\subsubsection{Professor Data}
$pt_{p}$ - Required number of sections to be taught by professor \\
$pq_{p,c}$ - Professor to course qualifications \\
$pcn_{p}$ - Indicates whether professor can teach at night or not \\
$pcd_{p}$ - Indicates whether professor can teach outside of DMS or not \\
$pcp_{p,c}$ - Professor preference for a course \\
$ptp_{p,s,td}$ - Professor preference for a time \\
$pbp_{p,bl}$ - Professor preference for a block \\
$ppp_{p}$ - Professor preference for a packed day 


\subsubsection{Classroom Data}
$kc_{k}$ - Classroom capacity\\
$\rho\in[0,1]$ - Minimum physical-room fill fraction, configured by
\texttt{minimum\_room\_fill}; the default is $0.8$


\subsubsection{Course \& Section Data}
$csa_{c,i,s}$ - Course section attendance \\
$cct_{c,c,s}$ - These courses $C$ must be scheduled together during $S$ \\
$ccn_{c,c,s}$ - These courses $C$ cannot be scheduled together during $S$ \\
$cp_{c}$ - Courses having part-timer penalty \\
$cf_{c}$ - Courses having friday night penalty \\
$cb_{c,bl}$ - Course block preference





%%%%%%%%%%%%%%%%
%%% DECISION VARIABLES
%%%%%%%%%%%%%%%%
\subsection{Decision Variables} 


\subsubsection{Main Decision Variables}
$X_{pcistk}$ -  assign a professor $p$ to a section $c,i$, semester $s$, time $t$, room $k$. 

A viable set of X, denoted $XV$, is defined by a following combination:
\begin{itemize}
	\item where $CS$ exists
	\item where $pq$ exists
	\item where physical-room attendance satisfies $\rho kc_k\le csa_{cis}\le kc_k$
	\item where 3 hour blocks with more than 75 students must be at night only
	\item where combinations of professors who can't teach at night and night blocks are removed
	\item where combinations of professors who can't teach outside of DMS and 'other' classrooms are removed
\end{itemize}

All declared professor--course qualifications are retained when constructing
$XV$. Course preferences rank eligible assignments; they do not remove
qualifications. Other eligibility rules, such as night-teaching permission,
still apply. The part-timer surrogate is exempt from named-faculty qualification
requirements. Placeholder assignments use the ``other'' room surrogate and do
not consume physical-room capacity.

In the formulation this restriction should generally be denoted in the summation over X, however it is placed outside of the equation for better readability\\ 

$X = \begin{cases}  integer, & \text{if } p \in \{\text{part-timer}\} \text{ or } t \in \{\text{placeholder} \} \text{ or } k \in \{\text{other}\}\\
binary, & \text{otherwise} \end{cases}$



\subsubsection{Penalty and Reward Decision Variables}
$TimeMax_{s}$ - Indicates a maximum of courses offered in each time of day, per semester $s$ \\
$TimeMin_{s}$ - Indicates a minimum of courses offered in each time of day, per semester $s$ \\ \\
$DayMax_{s}$ - Indicates a maximum of courses offered in each day of week, per semester $s$ \\
$DayMin_{s}$ - Indicates a minimum of courses offered in each day of week, per semester $s$ \\ \\
$KCount_{k}$ - Counts total sections assigned to each classrooms $k$ \\ 
$CBLPen_{c,bl}$ - Counts the number of violations of block preference $bl$ per course $c$   \\
$CFNPen_{c}$ - Counts the number of violations of blocks scheduled to Friday night 3 hour block \\
$PDayT_{p,s,dw}$ - Indicates if a professor $p$ is teaching during day of the week $dw$, semeseter $s$ \\ 
$PRoomU_{p,s,dw,k}$ - Indicates if a professor $p$ uses classroom $k$ during day of the week $dw$, semester $s$ \\
$PPhysical_{p,s,dw}$ - Indicates use of at least one physical room on that day\\
$RoomExtra_{p,s,dw}$ - Counts the classrooms a professor uses on a day past the first one \\
$ProfBBV_{}$ - Counts the number of back to back occurences for professor $p$ 


\subsubsection{Hard Constraint Violation Indicators}
$ConcI_{cst}$ - Counts the number of concurrencies of courses with itself \\
$ConcT_{c,cc,st}$ - Counts the number of concurrencies between courses where they should be together\\
$ConcN_{c,cc,st}$ - Counts the number of concurrencies between courses where they should not be together\\
$ConcO_{c,cc,s,t,tt}$ - Slack for separated courses assigned to distinct overlapping timeslots\\
$ProfTV{p}$ - Counts violations for required number of sections to be taught per professor


%%%%%%%%%%%%%%%%
%%% Constraints
%%%%%%%%%%%%%%%%
\subsection{Constraints}


\subsubsection{Professor Constraints}
Professors can't be scheduled multiple times during one time block
\begin{equation}
	\sum_{cik} X_{pcistk} \le 1 \quad \forall \quad P, S,T, 
	[pcistk] \in \{ XV \}, p \not \in {\text{part-timer}},  t \not \in \{\text{placeholder} \}
\end{equation}

Named professors cannot teach during overlapping real meeting times. All
summands below are restricted to $XV$:
\begin{equation}
	\sum_{cik} X_{pcistk} + \sum_{cik} X_{pcis,tt,k} \le 1
	\quad \forall p\in P^*,\ s\in S,\ [t,tt]\in tc.
\end{equation}

Professors must teach at least their required number of sections. The
nonnegative slack counts teaching above that load; it does not permit a shortfall.
\begin{equation}
	\sum_{cikst} X_{pcistk} = pt_{p} + ProfTV_{p} \quad \forall \quad p\in P^{*},\qquad ProfTV_p\ge0
\end{equation}

Professors can only teach courses they are qualified to teach. Unqualified
assignments are excluded by $XV$; qualification does not impose an additional
upper bound equal to the agreed teaching load.
\begin{equation}
	X_{pcistk}=0 \quad \text{if }p\in P^{*}\text{ and }pq_{pc}=0
\end{equation}

Counts the maximum number of days a professor is teaching per each semester
\begin{equation}
	\sum_{c,i,k,\, t \in T_{dw}} X_{pcistk} \le M \, PDayT_{p,s,dw} \quad \forall P, S, DW, [pcistk] \in \{ XV \}
\end{equation}

Counts the classrooms a professor uses per day, so room-hopping within a
teaching day can be penalised: the first room of the day is free, every
further room adds to $RoomExtra_{p,s,dw}$.
\begin{equation}
	\sum_{c,i,\, t \in T_{dw}} X_{pcistk} \le M \, PRoomU_{p,s,dw,k} 
	\quad \forall P, S, DW, K, [pcistk] \in \{ XV \}, k \not \in \{ \text{other} \}
\end{equation}
Let $K^*=K\setminus\{\text{other}\}$. The implementation uses a separate
physical-room-day indicator, since teaching only off campus must not subtract
a room from this count:
\begin{equation}
	PPhysical_{p,s,dw}\le\sum_{k\in K^*}PRoomU_{p,s,dw,k}
	\le |K^*|PPhysical_{p,s,dw}.
\end{equation}
The room-change expression is
$RoomExtra_{p,s,dw}=\sum_{k\in K^*}PRoomU_{p,s,dw,k}-PPhysical_{p,s,dw}$;
it is not a separate slack variable in the implementation.

Counts the number of occurences a professor has back to back classes (in different classrooms)
\begin{equation}
	\sum_{i} X_{pcistk} + X_{pci,s,tt,kk} = ProfBBV_{pcist,tt,k,kk}  \quad \forall \_ \_ \_ \text{ where }
	k \ne kk, [t, tt], \in \{ tb \}
\end{equation}

\subsubsection{Time of Day Constraints}
Counts the maximum and minimum number of sections offered over different times of day.
\begin{equation}
	\sum_{p,c,i,k,\, t \in T_{td}} X_{pcistk} \le TimeMax_{s} \quad \forall \quad S, TD, [pcistk] \in \{ XV \}
\end{equation}

\begin{equation}
	\sum_{p,c,i,k,\, t \in T_{td}} X_{pcistk} \ge TimeMin_{s} \quad \forall \quad S, TD, [pcistk] \in \{ XV \}
\end{equation}


\subsubsection{Day of Week Constraints}
Counts the maximum and minimum number of sections offered over different days of week. Since 1.5 hour blocks are distributed between 2 days, while 3 hour blocks distributed over 1 day. 

\begin{equation}
	\sum_{p,c,i,k,bl,\, t \in T_{dw} \cap T_{bl}} X_{pcistk} \, tbw_{bl} 
	\le DayMax_{s} \quad \forall \quad S, DW, [pcistk] \in \{ XV \}
\end{equation}

\begin{equation}
	\sum_{p,c,i,k,bl,\, t \in T_{dw} \cap T_{bl}} X_{pcistk} \, tbw_{bl} 
	\ge DayMin_{s} \quad \forall \quad S, DW, [pcistk] \in \{ XV \}
\end{equation}



\subsubsection{Classroom Constraints}
Each classroom can only be occupied once per period and the time slots that share periods must also be respected by classrooms.
\begin{equation}
	\sum_{pci} X_{pcistk} \le 1 \quad \forall \quad K,S,T,
	[pcistk] \in \{ XV \}, k \not \in \{ \text{other} \}, t \not \in \{\text{placeholder} \}
\end{equation}
\begin{equation}
	\sum_{pci} X_{pcistk} + \sum_{pci} X_{pcisttk} \le 1 \quad \forall \quad K,S,T,TT, 
	[t,tt] \in \{tc\}, [pcistk] \in \{ XV \}, k \not \in \{ \text{other} \}
\end{equation}

Physical-room eligibility requires attendance no greater than capacity and no
less than the configured fill fraction. Ineligible assignments are removed
from $XV$:
\begin{equation}
	X_{pcistk}=0 \quad\text{if } csa_{cis}<\rho kc_k
	\text{ or }csa_{cis}>kc_k,\qquad k\ne\text{other}.
\end{equation}
Setting $\rho=0$ disables only the minimum-fill restriction; capacity remains
an upper bound. The default $\rho=0.8$ preserves the institutional 80\% rule.

Finally the classroom utilization must be maximized which is done by counting the number of courses
allocated to classroom and that value is rewarded in the objective function. (this is not needed as it is modelled by minimizing "other" classrooms)
\begin{equation}
	\sum_{pcist} X_{pcistk} = KCount_{k} \quad \forall \quad K, 
	[pcistk] \in \{ XV \}, k \not \in \{ \text{other} \}
\end{equation}


\subsubsection{Course Constraints}
Every declared course--section--semester--time-of-day tuple is assigned exactly once,
including an explicit placeholder if no real meeting is selected.
\begin{equation}
	\sum_{p,\,t\in T_{td},\,k} X_{pcistk} = 1
	\quad \forall(c,i,s,td)\in CS.
\end{equation}

Three hour courses with more than 75 students can only be scheduled at night. (removed by $XV$)
\begin{equation}
	\sum_{pkt, TTD \ne 0, TBL = 1} X_{pcistk} = 0 \quad \forall \quad CS_{cis,td} \text{ where } csa_{cis} \ge 75, TD \not \in \{night \},  BL \in \{ 3hour \}
\end{equation}

Some courses are restricted to either 1.5 or 3 hour blocks.
\begin{equation}
	\sum_{p,i,s,k,\, t \in T_{bl}} X_{pcistk} = 0 + CBLPen_{c,bl} \quad \forall \quad C, BL, 
	[pcistk] \in \{ XV \}, [c,bl] \in \{ cb \}
\end{equation}

Some courses should not be scheduled on the 3 hour Friday night block.
\begin{equation}
	\sum_{p,i,s,k,\, t \in T_{\text{Fri}} \cap T_{\text{Night}}} X_{pcistk} = 0 + CFNPen_{c} 
	\quad \forall \quad C, [pcistk] \in \{ XV \}, c \in \{ cf \}
\end{equation}


\subsubsection{Concurrency Constraints}
The same-course penalty counts excess named-faculty sections at the same
meeting-time ID. This rule is unchanged by the distinct-time apart constraints
below; it does not itself add penalties across distinct overlapping IDs.
\begin{equation}
	\sum_{p \in P^{*},i,k} X_{pcistk} 
	\le 1 + ConcI_{cst} \quad \forall \quad C,S,T, 
	[pcistk] \in \{ XV \}, t \not \in \{\text{placeholder}\}
\end{equation}

Courses marked together require equal named-faculty assignment counts at each
exact meeting-time ID. The nonnegative penalty bounds the absolute difference;
merely overlapping meeting IDs do not satisfy this rule.
\begin{equation}
	\left|\sum_{p \in P^{*},i,k} X_{pcistk}
	- \sum_{p \in P^{*},i,k} X_{p,cc,istk}\right|
	\le ConcT_{c,cc,st}
	\quad [c,cc,s]\in cct,\quad t\text{ real}.
\end{equation}
The implementation expresses the absolute-value bound as two linear inequalities.

Some concurrent courses cannot be scheduled together ($[pcistk] \in \{ XV \}$ and $[p,cc,istk] \in \{ XV \}$)
\begin{equation}
	\sum_{p \in P^{*},i,k} X_{pcistk} 
	+ \sum_{p \in P^{*},i,k} X_{p,cc,istk} 
	\le 1  + ConcN_{c,cc,st} \quad \forall \quad C,CC,S,T, 
	t \not \in \{\text{placeholder}\}, [c,cc] \in \{ ccn \}
\end{equation}

For courses marked apart, distinct overlapping meeting IDs must also be
considered. Define the symmetric relation
$E=\{(t,tt),(tt,t):[t,tt]\in tc,\ t\ne tt,\ t,tt\text{ real}\}$.
For each distinct-course pair $[c,cc,s]\in ccn$ not marked together, impose
\begin{equation}
	\sum_{p\in P^{*},i,k}X_{pcistk}
	+\sum_{p\in P^{*},i,k}X_{p,cc,is,tt,k}
	\le 1+ConcO_{c,cc,s,t,tt}\qquad \forall(t,tt)\in E.
\end{equation}
Each orientation keeps the courses at opposite ends of the overlap pair.
No transitive closure is taken: overlap of $t$ with $u$ and $u$ with $v$ does
not imply overlap of $t$ with $v$. These are soft constraints with nonnegative
slack, like the same-time apart rule. Part-timer and placeholder assignments
remain exempt. The summed slack is an aggregate penalty, not necessarily a
count of distinct conflicting section pairs.


%%%%%%%%%%%%%%%%
%%% Objective Function
%%%%%%%%%%%%%%%%
\subsection{Objective Function}
The solver runs the following ten stages in order. Each of the first nine
minimizes a measure $f_j$ and adds the cap
$f_j(X)\le\lceil z_j-10^{-6}\rceil$, where $z_j$ is the returned incumbent
value. Thus later stages preserve the rounded cap, rather than re-optimizing
the preceding priority. The second stage caps its combined measure, not its
conflict and placeholder components separately.
\begin{enumerate}

	\item \textbf{Teaching load.}
		\begin{equation} 
			min \sum ProfTV_{p} \quad \Rightarrow \quad 
			\sum ProfTV_{p} \le \text{ value of } \sum ProfTV_{p}
		\end{equation}
		
	\item \textbf{Conflicts and unscheduled sections.}
		\begin{equation}
			\min f_2(X)=\sum ConcI+\sum ConcT+\sum ConcN+\sum ConcO
			+\sum_{t\in\{\text{placeholder}\}}X.
		\end{equation}

	\item \textbf{Off-campus assignments.} Count only assignments with real meeting
		times; unscheduled placeholders are already counted in stage 2.
		\begin{equation}
			\min f_3(X)=\sum_{\substack{p,c,i,s,\,t\text{ real}\\k=\text{other}}} X_{pcistk}
			\quad\Rightarrow\quad f_3(X)\le\lceil z_3-10^{-6}\rceil.
		\end{equation}

	\item \textbf{Part-timer room use.} Prefer physical rooms for named faculty.

	\item \textbf{Schedule balance.}
		\begin{equation}
		\begin{aligned}
		\min\quad & \sum_s (TimeMax_s-TimeMin_s)\\
		           & + \sum_s (DayMax_s-DayMin_s)
		\end{aligned}
		\end{equation}

	\item \textbf{Course block violations.} Minimize $\sum_{c,bl}CBLPen_{c,bl}$.

	\item \textbf{Friday-night penalties.} Minimize $\sum_c CFNPen_c$.

	\item \textbf{Room changes.} Minimize $\sum_{p,s,dw}RoomExtra_{p,s,dw}$.

	\item \textbf{Teaching days.} Minimize $\sum_{s,p,dw}PDayT_{p,s,dw}\,ppp_p$.

	\item \textbf{Preference rewards.} Maximize the sum of the following rewards:
			\begin{equation}
			\begin{aligned}
						 & + 1 \sum_{pcistk} X \, pcp_{pc} 		& [pcistk] \in XV, pcp_{pc} > 0 \\
						 & + 1 \sum_{pcistk,td} X \, ptp_{ps,td} & [pcistk] \in XV, [t,td] \in \{ TTD \}, ptp_{ps,td} > 0 \\
						 & + 1 \sum_{pcistk,bl} X 				& [pcistk] \in XV, [t,bl] \in \{ TBL \}, pc \in \{ pbp \} \\ 
			\end{aligned}
		\end{equation}
		
\end{enumerate}

Stages 1, 2, 4, 5, 6, 7, and 9 request a zero relative MIP gap. Stages 3,
8, and 10 allow a 5\% relative gap. A zero requested gap does not guarantee
that a time-limited run proves its optimum. Conversely, a solver's successful
termination with a positive requested gap does not establish exact mathematical
optimality. Exact lexicographic optimality requires proving every relevant
preceding optimum; this manuscript makes no such claim for a general run.

The configured time limit applies separately to each stage. Ten stages can
therefore consume approximately ten times that limit, plus model construction,
incumbent repair, and other overhead. Preference-definition rows are added
after stage 4. Their auxiliary values are repaired before offering the incumbent
to the next stage; room-use indicators are reconstructed again before stage 8.
These implementation choices have not been established here to improve runtime.

The following weighted breakdown is an explanatory score computed from the
final solution. It is not the single objective optimized by all ten stages;
only course, time, and block rewards enter the final maximization.

\begin{center}
\begin{tabular}{>{\raggedright\arraybackslash}p{0.38\textwidth}r p{0.42\textwidth}}
\toprule
\textbf{Component} & \textbf{Weight} & \textbf{What it measures} \\
\midrule
Time-of-day spread & $-10$ & Uneven distribution across times of day \\
Day-of-week spread & $-10$ & Uneven distribution across days \\
Course block violation & $-3$ & Assignment to a non-preferred block \\
Friday-night penalty & $-3$ & Penalized course assigned to Friday night \\
Room spread & $-3$ & Additional rooms used by one professor in one day \\
Course reward & $+1$ & Professor--course preference points \\
Time reward & $+1$ & Professor--time-of-day preference points \\
Block reward & $+1$ & Professor--block preference points \\
Packed-day penalty & $-1$ & Professor teaching on additional days \\
\bottomrule
\end{tabular}
\end{center}

\section{Scope, Limitations, and Evaluation Status}
This formulation describes an institutional faculty-allocation and course
scheduling application. Part-timer, unscheduled, and off-campus surrogates make
selected compromises explicit, but do not guarantee feasibility: named-faculty
load requirements and other hard constraints can still be incompatible.
The minimum-fill threshold and priority order are institutional policy choices,
not general claims about educational quality or fairness.

This manuscript reports no controlled benchmark experiments, independently
validated improvements over manual schedules, or cross-institution performance
results. Software regression tests check selected implementation properties;
they do not establish novelty, computational superiority, or publication
readiness. A research evaluation still needs a reconciled formulation and
independent schedule checker, documented instances, matched baselines, equal
total compute budgets, stage bounds and termination records, and reproducible
artifacts. Competition timetabling scores cannot be compared directly without
first matching the respective benchmark constraints and objectives.

\end{document}  
